\section{Conclusion: Receiving Councils Without Mythology}

The two Vatican Councils disclose a paradox of Catholic modernity. Vatican I strengthened the visible center as states, revolutions, and intellectual systems threatened ecclesial freedom and doctrinal unity. Its definition was narrower than its popular triumph, but the council's interruption left the center without an equally developed conciliar account of the whole. Vatican II used that strengthened center to recover episcopal, baptismal, scriptural, patristic, liturgical, and missionary dimensions and to authorize wide reform. The center then distributed change with a speed and uniformity that could undercut the organic and local realities the council praised.

The proper Catholic response is neither rupture nor amnesia. Twelve conclusions can now be stated.

\begin{enumerate}
\item \textbf{Vatican I taught more than infallibility.} \emph{Dei Filius}'s harmony of faith and reason and \emph{Pastor aeternus}'s primacy of jurisdiction are essential to the history.
\item \textbf{The ex cathedra dogma is true and tightly conditioned.} It is not personal omniscience, inspiration, impeccability, or a permanent state.
\item \textbf{The formula has interpretive liabilities.} Its unfinished ecclesial setting, negative consent clause, procedural under-specification, and popular vocabulary aided maximalism. Authoritative reception supplies the missing positive context without reversing the definition.
\item \textbf{Vatican II completes rather than cancels Vatican I.} The Church's infallibility, episcopal college, local Churches, whole people, and primatial head belong together.
\item \textbf{Vatican II made no new extraordinary dogmatic definition.} All sixteen documents were solemnly promulgated; their teachings carry differentiated authority and include reaffirmed dogma, authentic doctrine, discipline, and prudential judgment.
\item \textbf{Its silence about proper names did not reverse prior judgments.} \emph{Gaudium et spes} analyzes coercive systematic atheism without saying “Communism”; no text approves Masonry or religious indifferentism.
\item \textbf{The Council did not write Paul VI's Missal.} It mandated substantial reform and set limits. The Consilium, pope, Roman offices, conferences, translators, and local celebrants made successive choices.
\item \textbf{The reformed Missal is Catholic, valid, and lawful.} Those truths do not answer every historical question about scale, organic growth, options, translation, distribution, or pastoral fruit.
\item \textbf{The Western contraction is real.} Raw counts and per-Catholic rates show severe losses in religious life, priests, sacraments, attendance, and institutional density. Global growth and geographic transfer prevent universalizing the claim.
\item \textbf{Lay belief is fractured but not absent.} Survey instruments show both substantial Catholic instincts and widespread Eucharistic ignorance, dissent, and policy attitudes opposed to teaching. Wording and practice level materially change results.
\item \textbf{The causal model must be conjunctural.} Secularization, demography, sexual revolution, authorized revision, liturgical discontinuity, catechetical weakness, institute choices, dissent, and abuse interacted. Council and reform were material mediators, not sufficient causes.
\item \textbf{Schism is not reform.} Criticism and older worship can exist in full communion. Parallel bishops erected against express papal prohibition deny in act the primacy traditionalists claim to defend.
\end{enumerate}

The crisis should change how authority is exercised. Definitions need their ecclesial setting. Reform needs a demonstrated burden of necessity, memory, gradual reception, and honest measurement. Liturgy needs stability that precedes the celebrant. Catechesis must distinguish dogma, definitive doctrine, authentic teaching, discipline, and prudence. Bishops and popes must correct failed implementation without treating every correction as repudiation of predecessors.

The crisis should also change how tradition is loved. Tradition is neither the latest policy nor the year 1962 suspended forever. It is the living transmission of the apostolic deposit through Scripture, worship, doctrine, saints, pastors, and the whole Church across time. Its living character permits reform; its given character judges reform. The Roman primacy serves that continuity when it confirms the brethren, receives the Churches' witness, and refuses both state captivity and self-referential absolutism.

There is no statistical technique for resurrection. Numbers identify wounds and test claims; they do not regenerate faith. Renewal will require what both councils, properly received, actually teach: revealed truth, grace, sacrifice, holiness, reason, Scripture, Eucharistic worship, apostolic communion, missionary charity, and the courage to repent of institutional failure without fleeing the Church in which Christ placed the means of healing.

\begin{studyframe}
\textbf{Final answer.} Vatican I is the doctrinal and institutional precursor to Vatican II, not because its dogma caused a later revolution, but because its true yet compressed papal ecclesiology left a task and created a centralized instrument. Vatican II completed much of the task and used the instrument to authorize reform. Paul VI's Missal was the most pervasive postconciliar product of that authority. The subsequent crisis cannot be laid wholly at any one door; neither can ecclesial decisions be removed from its causes. Catholic fidelity is strong enough to affirm the dogmas, receive the council, celebrate a lawful Missal, recover neglected tradition, criticize failed prudence, and remain in communion at once.
\end{studyframe}
